\documentclass[11pt]{article}
\usepackage[utf8]{inputenc}
\usepackage[T1]{fontenc}
\usepackage[a4paper,margin=1in]{geometry}
\usepackage{lmodern}
\usepackage{microtype}
\usepackage{amsmath,amssymb,mathtools}
\usepackage{mathrsfs}
\usepackage{tikz-cd}
\usetikzlibrary{arrows}
\usepackage{enumitem}
\usepackage{longtable,booktabs,array}
\usepackage[hidelinks]{hyperref}
\setlength{\parindent}{0pt}
\setlength{\parskip}{0.55em}
\emergencystretch=3em
\hbadness=10000
\hfuzz=20pt
\usepackage[normalem]{ulem}
\newcommand{\K}{K}
\newcommand{\Gr}{\operatorname{Gr}}
\newcommand{\Fil}{\operatorname{Fil}}
\newcommand{\ch}{\operatorname{ch}}
\newcommand{\Todd}{\operatorname{Todd}}
\newcommand{\cl}{\operatorname{cl}}
\newcommand{\Parf}{\operatorname{Parf}}
\newcommand{\Tor}{\operatorname{Tor}}
\newcommand{\Spec}{\operatorname{Spec}}
\newcommand{\RRR}{\mathrm{RRR}}
\providecommand{\codim}{\operatorname{codim}}
\providecommand{\Inv}{\operatorname{Inv}}
\providecommand{\Pic}{\operatorname{Pic}}
\providecommand{\Z}{\mathbb Z}
\providecommand{\Q}{\mathbb Q}
\providecommand{\Ahat}{\widehat A}
\providecommand{\Ahp}{\widehat A^{+}}
\providecommand{\That}{\widehat T}
\providecommand{\rad}{\operatorname{rad}}
\providecommand{\Gl}{\operatorname{Gl}}
\providecommand{\GL}{\operatorname{Gl}}
\providecommand{\ch}{\operatorname{ch}}
\providecommand{\Todd}{\operatorname{Todd}}
\providecommand{\Gr}{\operatorname{Gr}}
\providecommand{\red}{\text{réd}}

% Core macro additions
\providecommand{\Atil}{\widetilde A}
\providecommand{\Ctil}{\widetilde C}
\providecommand{\K}{\mathrm K}
\providecommand{\cC}{\mathcal C}
\providecommand{\cF}{\mathscr F}
\providecommand{\eps}{\varepsilon}
\providecommand{\ellc}{\ell}
\providecommand{\Tor}{\operatorname{Tor}}
\providecommand{\Coker}{\operatorname{Coker}}
\providecommand{\Ker}{\operatorname{Ker}}
\providecommand{\Supp}{\operatorname{Supp}}
\providecommand{\codim}{\operatorname{codim}}
\providecommand{\Gr}{\operatorname{G}}
\providecommand{\rank}{\operatorname{rank}}
\providecommand{\rang}{\operatorname{rang}}
\providecommand{\cO}{\mathcal O}
\providecommand{\sgacircled}[1]{\tikz[baseline=(char.base)]\node[shape=circle,draw,inner sep=.35pt](char){\mbox{\scriptsize #1}};}

\hypersetup{hypertexnames=false}
\newcommand{\sgastandalonefootnotetext}[2]{%
  \begingroup
  \NoHyper
  \footnotetext{#2}%
  \endNoHyper
  \endgroup
}
\providecommand{\R}{\mathbb R}
\providecommand{\C}{\mathbb C}
\providecommand{\N}{\mathbb N}
\providecommand{\Zell}{\mathbb Z_\ell}
\providecommand{\Qell}{\mathbb Q_\ell}
\providecommand{\C}{\mathbb C}
\providecommand{\Ch}{\operatorname{Ch}}
\providecommand{\Sym}{\operatorname{Sym}}
\providecommand{\Td}{\operatorname{Td}}
\begin{document}
% Additional macro definitions
\providecommand{\cC}{\mathcal C}
\providecommand{\cD}{\mathcal D}
\providecommand{\cK}{\mathcal K}
\providecommand{\cO}{\mathcal O}
\providecommand{\cS}{\mathcal S}
\providecommand{\Cpx}{\mathrm C}
\providecommand{\Kpx}{\mathrm K}
\providecommand{\Dpx}{\mathrm D}
\providecommand{\D}{\mathrm D}
\providecommand{\Ho}{\mathrm H}
\providecommand{\RHom}{\mathrm R\!\operatorname{Hom}}
\providecommand{\Ob}{\operatorname{ob}}
\providecommand{\Coker}{\operatorname{Coker}}
\providecommand{\Ker}{\operatorname{Ker}}
\renewcommand{\Im}{\operatorname{Im}}
\providecommand{\Hom}{\operatorname{Hom}}
\providecommand{\Cone}{\operatorname{Cone}}
\providecommand{\id}{\operatorname{Id}}
\providecommand{\qis}{\xrightarrow{\sim}}


%%%%%%%%%%  scan idx 98  |  folio 85  %%%%%%%%%%

% Additional notation for Expose I continuation
\providecommand{\RHom}{\mathrm R\!\Hom}
\providecommand{\Ext}{\operatorname{Ext}}
\providecommand{\coh}{\operatorname{coh}}
\providecommand{\PreCoh}{\operatorname{precoh}}
\providecommand{\Tr}{\operatorname{Tr}}
\providecommand{\Qis}{\operatorname{Qis}}
\providecommand{\Rad}{\operatorname{rad}}
\providecommand{\Spec}{\operatorname{Spec}}
\providecommand{\End}{\operatorname{End}}
\providecommand{\Id}{\operatorname{Id}}
\providecommand{\varinjlim}{\mathop{\mathrm{lim}}\limits_{\longrightarrow}}
\providecommand{\varprojlim}{\mathop{\mathrm{lim}}\limits_{\longleftarrow}}
\providecommand{\qcoh}{\operatorname{qcoh}}
\providecommand{\amp}{\operatorname{amp}}

%%%%%%%%%%  scan idx 133  |  folio 120  %%%%%%%%%%
\emph{un
diagramme de \(\Cpx(\mathcal C_S)\), où \(P\) est strictement \(\mathcal C_0\)-parfait et \(f\) un quasi-isomorphisme. L'ensemble des \(X\) au-dessus de \(S\) tels qu'il existe une flèche \(u:P_X\to E_X\) de \(\Cpx(\mathcal C_X)\) telle que le diagramme}
\[
\begin{tikzcd}
& P_X \arrow[d,"v_X"] \arrow[dl,"u"']\\
E_X \arrow[r,"f_X"] & F_X
\end{tikzcd}
\]
\emph{soit commutatif dans \(\Kpx(\mathcal C_X)\) est un raffinement de \(S\).}

\emph{Preuve.} Soit \(G\) le cône de \(f\). On a la suite exacte
\[
  \longrightarrow \Hom_{\Kpx(\mathcal C_S)}(P,E)
  \longrightarrow \Hom_{\Kpx(\mathcal C_S)}(P,F)
  \longrightarrow \Hom_{\Kpx(\mathcal C_S)}(P,G)
  \longrightarrow .
\]
Comme \(G\) est acyclique, l'image de \(v\) dans \(\Hom_{\Kpx(\mathcal C_S)}(P,G)\) est localement nulle d'après (4.1). Donc, quitte à localiser, il existe \(u:P\to E\) tel que \(v=fu\) dans \(\Kpx(\mathcal C_S)\), cqfd.

\textbf{Corollaire 4.3.} \emph{Soit \(f:E\to F\) un quasi-isomorphisme de \(\Cpx(\mathcal C_S)\), avec \(F\) strictement parfait (relativement à \(\mathcal C_0\)). L'ensemble des \(X\) au-dessus de \(S\) tels qu'il existe un quasi-isomorphisme \(g:F_X\to E_X\) tel que \(f_Xg\) soit homotope à \(\Id_{F_X}\) est un raffinement de \(S\).}

\emph{Preuve.} C'est un cas particulier de (4.2).

\textbf{Scholie 4.4.} Soit \(F\in\Ob\Kpx(\mathcal C_S)\). On peut définir (par exemple à l'aide d'un clivage de \(\mathcal C\) (SGA 1 VI 7.1)) une catégorie \((\underline{\Qis/F})\) fibrée sur \(S_{/S}\), dont la fibre en \(X\) au-dessus de \(S\) est la catégorie \((\Qis/F_X)^\circ\) (opposée de la catégorie des quasi-isomorphismes de but \(F_X\)). On peut alors exprimer (4.3) en disant que, 
%%%%%%%%%%  scan idx 134  |  folio 121  %%%%%%%%%%
si \(F\) est
strictement parfait, la flèche \(\Id_F\) est localement cofinale dans \((\underline{\Qis/F})\).

\textbf{Notation 4.5.} Soit \(\mathcal M\) une \(S\)-catégorie fibrée. Soient \(X\in\Ob S\), et \(E,F\in\Ob\mathcal M_X\). Nous noterons \(\underline{\Hom}_{\mathcal M}(E,F)\) le faisceau sur \(S_{/X}\) associé au préfaisceau des \(S\)-morphismes de \(E\) dans \(F\) dans \(\mathcal M\) (pour la définition de ce préfaisceau, voir [G] I 2.6; on prendra garde que, même si \(\mathcal C\) est un champ ([G]), le préfaisceau correspondant pour \(\mathcal M=\Kpx(\mathcal C)\) (ou \(\Dpx(\mathcal C)\)) n'est pas en général un faisceau).

\textbf{Corollaire 4.6.} \emph{Soient \(E,F\in\Ob\Cpx(\mathcal C_S)\), \(E\) étant strictement parfait. La flèche canonique}
\[
  \underline{\Hom}_{\Kpx(\mathcal C)}(E,F)
  \longrightarrow
  \underline{\Hom}_{\Dpx(\mathcal C)}(E,F)
\]
\emph{est un isomorphisme}.

\emph{Preuve.} C'est une conséquence immédiate de (4.4) et de la formule
\[
  \Hom_{\Dpx(\mathcal C_X)}(E_X,F_X)
  = \varinjlim_{(\Qis/E_X)^\circ}\Hom_{\Kpx(\mathcal C_X)}(\,\cdot\,,F_X)
\]
pour \(X\) au-dessus de \(S\).

\textbf{Définition 4.7.} \emph{Soit \(F\in\Ob\Dpx(\mathcal C_S)\) et soient \(a,b\in\mathbb Z\) tels que \(a\leq b\). Nous dirons que \(F\) est d'}amplitude \(\mathcal C_0\)-parfaite contenue dans l'intervalle \([a,b]\), \emph{et nous écrirons}
\[
  \mathcal C_0\text{-}\mathrm{parf}\text{-}\amp(F)\subset[a,b],
\]
\emph{si l'ensemble des \(X\) au-dessus de \(S\) tels qu'il existe un quasi-isomorphisme} 
%%%%%%%%%%  scan idx 135  |  folio 122  %%%%%%%%%%
\(F'\to F\) \emph{de} \(\Cpx(\mathcal C_X)\)
\emph{avec \(F'\) strictement \(\mathcal C_0\)-parfait (2.1), et tel que \(F'^i=0\) pour \(i\notin[a,b]\), est un raffinement de} \(S\). (\emph{En vertu de (4.3) cette condition équivaut à celle obtenue en remplaçant l'expression « quasi-isomorphisme de \(\Cpx(\mathcal C_X)\) » par « isomorphisme de \(\Dpx(\mathcal C_X)\) ».}) \emph{Nous dirons que \(F\) est d'}amplitude \(\mathcal C_0\)-parfaite finie \emph{s'il existe des entiers \(a\) et \(b\) tels que} \(\mathcal C_0\)-\(\mathrm{parf}\)-\(\amp(F)\subset[a,b]\). \emph{Nous dirons que \(F\) est} \(\mathcal C_0\)-parfait \emph{si \(F\) est localement d'amplitude \(\mathcal C_0\)-parfaite finie}\footnote{Bien entendu, quand cela ne pourra prêter à confusion, nous omettrons \(\mathcal C_0\) dans les expressions précédentes.}.

\textbf{Exemple 4.8.} Soient \((S,A)\) un topos annelé, \(\mathcal C\) la \(S\)-catégorie définie par \(U\longmapsto\) (catégorie des \(A_U\)-Modules à gauche), \(\mathcal C_0\) la sous-\(S\)-catégorie définie par \(U\longmapsto\) (catégorie des \(A_U\)-Modules à gauche facteurs directs de Modules libres de type fini). Le couple \((\mathcal C,\mathcal C_0)\) satisfait, d'après (1.3.1), aux hypothèses (4.0). Un objet \(F\) de \(\Dpx(_A S)\) (cf. (2.15)) sera dit d'amplitude parfaite contenue dans \([a,b]\) (resp. d'amplitude parfaite finie, resp. parfait) si \(F\) est d'amplitude \(\mathcal C_0\)-parfaite contenue dans \([a,b]\) (resp.\dots). Lorsque les facteurs directs de Modules libres de type fini sont localement libres (cf. (2.15.1)), dire que \(F\) est parfait signifie donc que \(F\) est localement isomorphe (dans \(\Dpx(S)\)) à un complexe à degrés bornés et à composantes libres de type fini.

\textbf{Notation 4.9.} Nous désignerons par \(\Dpx(\mathcal C)_{\mathcal C_0\text{-}\mathrm{parf}}\) ou \(\Parf(\mathcal C,\mathcal C_0)\) la sous-\(S\)-catégorie strictement pleine de \(\Dpx(\mathcal C)\) formée des objets \(\mathcal C_0\)-parfaits. Dans le cas de l'exemple (4.8), nous écrirons \(\Dpx(_A S)_{\mathrm{parf}}\), ou \(\Parf(_A S)\), ou même \(\Parf(S)\) s'il n'y a pas d'ambiguïté possible sur \(A\).

%%%%%%%%%%  scan idx 136  |  folio 123  %%%%%%%%%%
\textbf{Proposition 4.10} \emph{La catégorie \(\Parf(\mathcal C,\mathcal C_0)\) est une sous-\(S\)-catégorie triangulée de \(\Dpx(\mathcal C)\) (1.1.4). Le foncteur canonique}
\[
  \Kpx^b(\mathcal C_0)\longrightarrow\Parf(\mathcal C,\mathcal C_0)
\]
(resp. \(\Tr(\Kpx^b(\mathcal C_0))\to\Tr(\Parf(\mathcal C,\mathcal C_0))\) (1.1.3)) \emph{induit une \(S\)-équivalence sur les champs associés ([G] II 2.2.4), ce qui signifie} :
\begin{enumerate}[label=(\roman*)]
\item \emph{pour tout \(X\in\Ob S\) et tous} \(E,F\in\Ob\Kpx^b(\mathcal C_{0X})\) (resp. \(\Tr(\Kpx^b(\mathcal C_{0X}))\)), \emph{la flèche canonique}
\[
  \underline{\Hom}_{\Kpx^b(\mathcal C_0)}(E,F)
  \longrightarrow
  \underline{\Hom}_{\Dpx(\mathcal C)}(E,F)
\]
(\emph{resp.} \(\underline{\Hom}_{\Tr(\Kpx^b(\mathcal C_0))}(E,F)\to\underline{\Hom}_{\Tr(\Dpx(\mathcal C))}(E,F)\)) \emph{est un isomorphisme, et}
\item \emph{tout objet de \(\Parf(\mathcal C,\mathcal C_0)\) (resp. \(\Tr(\Parf(\mathcal C,\mathcal C_0))\)) est localement isomorphe dans \(\Dpx(\mathcal C)\) (resp.} \(\Tr(\Dpx(\mathcal C))\) \emph{à un objet de \(\Kpx^b(\mathcal C_0)\) (resp.} \(\Tr(\Kpx^b(\mathcal C_0))\).
\end{enumerate}

\emph{Preuve.} L'assertion (i) non respée est un cas particulier de (4.6) et l'assertion (ii) non respée est vraie par définition ((4.7) et (4.9)). Il est clair d'autre part que \(\Parf(\mathcal C,\mathcal C_0)\) est stable par translation des degrés. Pour prouver que \(\Parf(\mathcal C,\mathcal C_0)\) est une sous-\(S\)-catégorie triangulée de \(\Dpx(\mathcal C)\), il suffit donc de montrer que si
\[
\begin{tikzcd}
& G \arrow[dl,"+1"'] & \\
E \arrow[rr,"u"] && F \arrow[ul]
\end{tikzcd}
\tag{*}
\]
est un triangle distingué de \(\Dpx(\mathcal C)\) tel que \(E\) et \(F\) soient parfaits, alors \(G\) 
%%%%%%%%%%  scan idx 137  |  folio 124  %%%%%%%%%%
est
parfait. Quitte à localiser, on peut supposer que \(E\) et \(F\) sont objets de \(\Kpx^b(\mathcal C_0)\). Quitte à localiser davantage, on peut supposer, d'après l'assertion (i) non respée, que \(u\) est une flèche de \(\Kpx^b(\mathcal C_0)\); \(G\) est alors parfait car isomorphe au cône de \(u\), qui est strictement parfait. Donc \(\Parf(\mathcal C,\mathcal C_0)\) est bien une sous-\(S\)-catégorie triangulée de \(\Dpx(\mathcal C)\). L'argument précédent prouve aussi, d'ailleurs, l'assertion (ii) respée. Reste à prouver l'assertion (i) respée. Soient donc \(X\in\Ob S\), et \(E=(E_1\xrightarrow{a_1}E_2\xrightarrow{a_2}E_3\xrightarrow[+1]{a_3}E_1)\), \(F=(F_1\xrightarrow{b_1}F_2\xrightarrow{b_2}F_3\xrightarrow[+1]{b_3}F_1)\) des triangles distingués de \(\Kpx^b(\mathcal C_{0X})\), montrons que la flèche
\[
  \underline{\Hom}_{\Tr(\Kpx^b(\mathcal C_0))}(E,F)
  \longrightarrow
  \underline{\Hom}_{\Tr(\Dpx(\mathcal C))}(E,F)
  \tag{**}
\]
est un isomorphisme. D'après l'assertion (i) non respée, on sait déjà que \((**)\) est injective. Soit \((u_i):E\to F\) un morphisme de triangles dans \(\Dpx(\mathcal C)\). Quitte à localiser, on peut supposer, d'après l'assertion (i) non respée, que \(u_i\) (\(1\leq i\leq3\)) est une flèche de \(\Kpx^b(\mathcal C_0)\). La flèche \(u_{i+1}a_i-b_iu_i\), pour \(1\leq i\leq3\) (avec la convention \(u_4=u_1\)), est nulle dans \(\Dpx(\mathcal C)\), donc localement nulle dans \(\Kpx^b(\mathcal C_0)\) d'après l'assertion (i) non respée. En d'autres termes, \((u_i)\) définit localement un morphisme de triangles dans \(\Kpx^b(\mathcal C_0)\). Cela prouve que \((**)\) est surjective et achève la démonstration de (4.10).

\textbf{Remarque 4.10.1.} On peut remplacer, dans (4.10), \(\Parf(\mathcal C,\mathcal C_0)\) par \(\Dpx^*(\mathcal C)_{\mathcal C_0\text{-}\mathrm{parf}}=\Dpx^*(\mathcal C)\cap\Parf(\mathcal C,\mathcal C_0)\) \((*=-\text{ ou }b)\), il n'y a rien à changer à la démonstration.

\textbf{Complément 4.11.} On peut préciser la structure triangulée de \(\Parf(\mathcal C,\mathcal C_0)\) en termes d'amplitude parfaite (4.7). Si \(E\in\Ob\Dpx(\mathcal C)\) est d'amplitude parfaite contenue dans \([a,b]\), \(E[n]\) est d'amplitude parfaite contenue dans \([a-n,b-n]\).

%%%%%%%%%%  scan idx 138  |  folio 125  %%%%%%%%%%
Soit
\[
\begin{tikzcd}
& G \arrow[dl,"+1"'] & \\
E \arrow[rr] && F \arrow[ul]
\end{tikzcd}
\]
un
triangle distingué de \(\Dpx(\mathcal C)\). On peut dresser un tableau analogue à (2.6), où les cases indiquent des amplitudes parfaites et chaque ligne correspond à une implication dont le but est souligné.
\[
\begin{array}{ccc}
E&F&G\\[0.3em]
{}[a+1,b+1]&{}[a,b]&\underline{[a,b]}\\
{}[a,b]&\underline{[a,b]}&{}[a,b]\\
\underline{[a,b]}&{}[a,b]&{}[a-1,b-1].
\end{array}
\]
De (4.10) on déduit aussitôt le

\textbf{Corollaire 4.12.} \emph{Supposons que \(S\) soit un site ponctuel. Pour que \(F\in\Ob\Dpx(\mathcal C)\) soit d'amplitude parfaite (rel.\ à \(\mathcal C_0\)) contenue dans \([a,b]\), il faut et il suffit qu'il existe un quasi-isomorphisme \(F'\to F\), avec \(F'\) strictement parfait et tel que \(F'^i=0\) pour \(i\notin[a,b]\). Les foncteurs canoniques}
\[
  \Kpx^b(\mathcal C_0)\longrightarrow\Parf(\mathcal C,\mathcal C_0)
\]
\emph{et}
\[
  \Tr(\Kpx^b(\mathcal C_0))\longrightarrow\Tr(\Parf(\mathcal C,\mathcal C_0))
\]
\emph{sont des équivalences de catégories}.

\textbf{Remarque 4.12.1.} Quand \(S\) est un site ponctuel, les objets de \(\mathcal C_0\) sont projectifs. Lorsque \(\mathcal C_0\) contient tous les objets projectifs de \(\mathcal C\), on dit « amplitude projective » au lieu de « amplitude \(\mathcal C_0\)-parfaite ».



% Additional notation for Expose I continuation
\providecommand{\Cpx}{\mathrm C}
\providecommand{\Kpx}{\mathrm K}
\providecommand{\Dpx}{\mathrm D}
\providecommand{\Ho}{\mathrm H}
\providecommand{\RHom}{\mathrm R\!\operatorname{Hom}}
\renewcommand{\Ob}{\operatorname{ob}}
\providecommand{\qis}{\xrightarrow{\sim}}
\providecommand{\RHom}{\mathrm R\!\Hom}
\providecommand{\Ext}{\operatorname{Ext}}
\providecommand{\Hom}{\operatorname{Hom}}
\providecommand{\Tr}{\operatorname{Tr}}

\providecommand{\Add}{\operatorname{Add}}
\providecommand{\Cart}{\underline{\operatorname{Cart}}}
\providecommand{\perfamp}{\operatorname{parf.amp}}
\providecommand{\perfdim}{\operatorname{parf.dim}}
\providecommand{\toramp}{\operatorname{tor.amp}}
\providecommand{\torf}{\operatorname{torf}}
\providecommand{\injf}{\operatorname{injf}}
\providecommand{\rg}{\operatorname{rg}}
\providecommand{\is}{\mathrm{is}}
\providecommand{\Mod}{\operatorname{Mod}}



\end{document}
